\chapter{Anxiety}

\section*{Definition and Overview}
Anxiety is a normal and useful emotion -- a feeling of fear, worry, or apprehension that alerts us to danger and mobilises the body to respond. It becomes an \textbf{anxiety disorder} when the fear is excessive, persistent, or out of proportion to the situation, and begins to interfere with daily life. Anxiety disorders are a group of related conditions, including generalised anxiety disorder, panic disorder, specific phobias, social anxiety disorder, agoraphobia, and separation anxiety disorder.

They are among the most common mental health conditions in the world, and, importantly, they are highly treatable: most people improve significantly with psychological therapy, medication, or both. Anxiety commonly coexists with depression, and the two conditions share many of the same treatments.

\section*{Epidemiology}
Anxiety disorders affect roughly one in four adults at some point in their lives, and around one in ten in any given year. They are more common in women than in men, and frequently begin in childhood, adolescence, or early adulthood. Generalised anxiety disorder and social anxiety are among the most prevalent forms.

Anxiety disorders often occur alongside depression, substance misuse, irritable bowel syndrome, and chronic pain, and they carry substantial personal and economic costs through reduced work productivity and healthcare use. Many people with anxiety do not seek help, so the condition is often under-diagnosed. Rates have been rising among young people, partly attributed to academic pressure, social media, and greater awareness and reporting.

\section*{Aetiology and Pathophysiology}
Anxiety disorders arise from an interaction of biological, psychological, and social factors rather than any single cause:

\begin{itemize}
    \item \textbf{Genetics:} anxiety runs in families, and twin studies indicate a substantial inherited component
    \item \textbf{Brain function:} overactivity of the amygdala (the brain's fear centre) and altered regulation of the amygdala by the prefrontal cortex, together with disruption of the HPA stress axis
    \item \textbf{Neurotransmitters:} imbalances or altered sensitivity involving serotonin, noradrenaline, and gamma-aminobutyric acid (GABA)
    \item \textbf{Temperament:} people who are naturally cautious, sensitive to stress, or prone to negative thinking are more vulnerable
    \item \textbf{Life events:} childhood trauma, abuse, major losses, and chronic stress can trigger or worsen anxiety
    \item \textbf{Learned responses:} avoidance of feared situations reinforces the fear and maintains the disorder
    \item \textbf{Substances:} caffeine, alcohol, stimulants, and some prescribed medicines can provoke or mimic anxiety
\end{itemize}

\section*{Clinical Features}
\begin{itemize}
    \item Persistent, hard-to-control worry or dread that is out of proportion to the situation
    \item Restlessness, feeling constantly on edge, or irritability
    \item Poor concentration, racing thoughts, and difficulty making decisions
    \item Sleep problems, including trouble falling or staying asleep and waking unrefreshed
    \item Muscle tension, headaches, trembling, and fatigue
    \item Physical symptoms such as palpitations, sweating, dizziness, breathlessness, and stomach upset
    \item Panic attacks: sudden surges of intense fear with a racing heart, chest tightness, and a sense of losing control or dying
    \item Avoidance of feared situations, places, or social events, which can progressively narrow a person's life
\end{itemize}

\section*{Differential Diagnosis}
\begin{itemize}
    \item \textbf{Depression:} which often coexists with anxiety and can produce similar physical and cognitive symptoms
    \item \textbf{Bipolar disorder} and other mental health conditions with agitation or worry
    \item \textbf{Attention deficit hyperactivity disorder:} restlessness and poor concentration can resemble anxiety
    \item \textbf{Post-traumatic stress disorder:} fear related to a specific trauma
    \item \textbf{Hyperthyroidism and other hormonal disorders:} which cause palpitations, sweating, and tremor
    \item \textbf{Cardiac conditions:} arrhythmias, mitral valve disease, or ischaemia causing palpitations and breathlessness
    \item \textbf{Respiratory conditions:} asthma and hyperventilation syndromes
    \item \textbf{Side effects of medicines:} such as stimulants, thyroid hormones, and some decongestants
    \item \textbf{Substance misuse or withdrawal:} particularly alcohol, benzodiazepines, and stimulants
    \item \textbf{Caffeine excess:} a common and easily overlooked cause of anxiety symptoms
\end{itemize}

\section*{When to Seek Medical Care}
See a GP if anxiety is persistent, severe, or stopping you from working, studying, or socialising. It is worth seeking help even if symptoms have been present for years, because effective treatments are available. Tell your doctor about any new or worrying chest pain, palpitations, or breathlessness so that physical causes can be properly excluded.

\textbf{Call 000 (or local emergency services) for:}
\begin{itemize}
    \item Thoughts of suicide or self-harm, or plans to end your life
    \item A panic attack with severe chest pain, fainting, or difficulty breathing
    \item Feeling unable to keep yourself safe or care for yourself
    \item New confusion, or a rapid, irregular heartbeat
\end{itemize}

\textbf{Immediate support} is also available from Lifeline (13 11 14) and other crisis services.

\section*{Diagnosis and Investigations}
\begin{itemize}
    \item \textbf{Clinical interview:} a detailed discussion of symptoms, their onset and triggers, and their impact on work, study, relationships, and sleep
    \item \textbf{Screening questionnaires:} tools such as the GAD-7 to measure severity, monitor progress, and guide treatment decisions
    \item \textbf{Physical examination:} to check for conditions that can cause anxiety symptoms, including heart rate, blood pressure, and thyroid palpation
    \item \textbf{Blood tests:} thyroid function, full blood count, and blood glucose where medical causes are suspected
    \item \textbf{Substance and medication review:} assessment of caffeine, alcohol, recreational drugs, and prescribed medicines
    \item \textbf{Co-morbidity screen:} structured assessment for depression and other mental health conditions
\end{itemize}

\section*{Management}
\begin{itemize}
    \item \textbf{Cognitive behaviour therapy (CBT):} the most effective psychological treatment, teaching people to challenge unhelpful thoughts and reduce avoidance
    \item \textbf{Exposure therapy:} gradually and safely facing feared situations under guidance, which reduces fear over time
    \item \textbf{Relaxation and mindfulness:} breathing exercises, progressive muscle relaxation, and mindfulness-based approaches for panic and worry
    \item \textbf{Medication:} SSRIs are the first-line drug treatment; short courses of benzodiazepines may be used briefly in specific circumstances, always under medical supervision
    \item \textbf{Lifestyle measures:} regular exercise, consistent sleep, and limiting caffeine and alcohol
    \item \textbf{Structured support:} a GP Mental Health Treatment Plan, referral to a psychologist or psychiatrist, and peer support groups
    \item \textbf{Treating co-morbidity:} addressing depression, substance use, and physical health conditions improves overall outcome
\end{itemize}

\section*{Complications}
\begin{itemize}
    \item Depression and an increased risk of suicide
    \item Substance misuse, most often alcohol, as a form of self-medication
    \item Avoidance leading to social isolation and lost work or study opportunities
    \item Strain on relationships and family life
    \item Physical consequences of chronic stress, including high blood pressure, heart disease, and disordered sleep
    \item Reduced quality of life and, in some people, chronic disability
    \item Impaired driving or workplace safety during panic episodes
\end{itemize}

\section*{Prognosis and Outcomes}
Most people with anxiety disorders improve substantially with treatment. CBT helps many people manage their symptoms well, and medication can reduce symptoms while therapy takes effect; the combination is often superior to either alone. Some people recover completely, while others learn to live well with a tendency toward anxiety that may return during times of stress.

Without treatment, anxiety disorders typically follow a long, fluctuating course, which is why early and consistent treatment is encouraged. Good prognostic factors include early treatment, adherence to therapy, social support, and the absence of severe co-morbid depression or substance misuse.

\section*{Prevention and Self-Care}
\begin{itemize}
    \item Learn relaxation and controlled-breathing techniques and practise them regularly, not only in a crisis
    \item Keep up regular physical activity -- it reduces anxiety and improves mood
    \item Maintain a regular sleep routine and address sleep problems early
    \item Limit caffeine and alcohol, and avoid stimulant drugs
    \item Notice early warning signs and seek help before anxiety becomes severe
    \item Face feared situations gradually rather than avoiding them
    \item Stay connected with family and friends, and talk about how you feel
    \item Set realistic goals and allow time for rest and enjoyment
\end{itemize}

\section*{Sources and Further Reading}
The following organisations provide reliable, up-to-date information on anxiety and its treatment for patients, students, and health professionals.

\begin{itemize}
    \item NHS. \textit{Information on anxiety disorders, symptoms, and treatment.} https://www.nhs.uk/conditions/
    \item NICE. \textit{Guidelines on generalised anxiety and panic disorder in adults.} https://www.nice.org.uk/
    \item Mayo Clinic. \textit{Overview of anxiety disorders, symptoms, and treatment.} https://www.mayoclinic.org/
    \item MedlinePlus. \textit{Health information on anxiety disorders and panic disorder.} https://medlineplus.gov/
    \item World Health Organization. \textit{Resources on mental health and anxiety.} https://www.who.int/
    \item MSD Manual. \textit{Professional overview of anxiety disorders.} https://www.msdmanuals.com/
\end{itemize}
